\documentclass[dvipsnames]{book}

\usepackage[margin=1cm,a5paper]{geometry}

\usepackage{bold-extra}
\usepackage{polski}
\usepackage{iwona}
\usepackage[OT4]{fontenc}

\usepackage{amsmath,amssymb,amsthm}
\usepackage{pgf,tikz}
\usepackage{hyperref}

\pagestyle{empty}
\setlength{\parindent}{0pt}
\setlength{\parskip}{5pt}
\renewcommand\leq\leqslant
\renewcommand\geq\geqslant

\colorlet{maincolor}{ForestGreen}
\usepackage[most]{tcolorbox}
\usepackage{varwidth}

\newtcolorbox{tresc}[3][]{enhanced, sharp corners, boxrule=0mm, colback=maincolor!10,
colframe=black,toptitle=1mm,fonttitle=\bfseries,
colbacktitle=maincolor!25, coltitle=black,
title={Zadanie #2\hfill{\it Opracowanie: #3}},#1}

\newenvironment{problem}[3]{\newpage\begin{tresc}{#1}{#2}{#3}}{\end{tresc}}

% Dodatkowe makra

\usepackage{multicol}
\usepackage{cancel}
\usepackage{bbm}

\def\ind{\mathbbm{1}}

\def\N{\mathbb{N}}
\def\P{\mathbb{P}}
\def\E{\mathbb{E}}
\def\Z{\mathbb{Z}}
\def\Q{\mathbb{Q}}
\def\R{\mathbb{R}}

\def\nwd{\mathrm{NWD}}

\newtheorem*{lemat}{Lemat}

\usepackage{enumitem}
\setlist[enumerate,1]{label={(\alph*)},topsep=0pt}

% ------------------------------------------------------------------
% Kolorowe boxy: definicje (zielone), twierdzenia/obserwacje/wnioski (niebieskie)
% ------------------------------------------------------------------

\colorlet{defcolor}{ForestGreen}
\colorlet{thmcolor}{RoyalBlue}

\newcounter{defcnt}
\newcounter{thmcnt}

% wspólny styl boxa
\tcbset{notatkibox/.style={enhanced, sharp corners, boxrule=0mm, frame hidden,
  toptitle=0.5mm, bottomtitle=0.5mm, left=2mm, right=2mm, top=1.5mm, bottom=1.5mm,
  fonttitle=\bfseries, coltitle=black, breakable, before skip=6pt, after skip=6pt}}

% \begin{definicja}[nazwa]
\NewDocumentEnvironment{definicja}{o}{%
  \refstepcounter{defcnt}%
  \begin{tcolorbox}[notatkibox, colback=defcolor!8, colframe=defcolor!25,
    colbacktitle=defcolor!25,
    title={Definicja \thedefcnt\IfValueT{#1}{\ (#1)}}]%
}{\end{tcolorbox}}

% generator środowisk niebieskich: #1 = nazwa środowiska, #2 = etykieta
\newcommand{\newthmbox}[2]{%
  \NewDocumentEnvironment{#1}{o}{%
    \refstepcounter{thmcnt}%
    \begin{tcolorbox}[notatkibox, colback=thmcolor!8, colframe=thmcolor!25,
      colbacktitle=thmcolor!25,
      title={#2 \thethmcnt\IfValueT{##1}{\ (##1)}}]%
  }{\end{tcolorbox}}%
}

\newthmbox{twierdzenie}{Twierdzenie}
\newthmbox{obserwacja}{Obserwacja}
\newthmbox{wniosek}{Wniosek}
\newthmbox{stwierdzenie}{Stwierdzenie}
% uwaga -- żółty box, \begin{uwaga}[nazwa]
\colorlet{uwagacolor}{Dandelion}

\NewDocumentEnvironment{uwaga}{o}{%
  \refstepcounter{thmcnt}%
  \begin{tcolorbox}[notatkibox, colback=uwagacolor!15, colframe=uwagacolor!40,
    colbacktitle=uwagacolor!40,
    title={Uwaga \thethmcnt\IfValueT{#1}{\ (#1)}}]%
}{\end{tcolorbox}}

% treść zadania -- szary box, \begin{zadanie}{1.1}
\colorlet{zadcolor}{Gray}

\NewDocumentEnvironment{zadanie}{m}{%
  \begin{tcolorbox}[notatkibox, colback=zadcolor!15, colframe=zadcolor!40,
    colbacktitle=zadcolor!40,
    title={Zadanie #1}]%
}{\end{tcolorbox}}

% przykład -- bez boxa, wyróżniony nagłówkiem
\theoremstyle{definition}
\newtheorem*{przyklad}{Przykład}

\renewcommand{\qedsymbol}{$\square$}


\begin{document}

\begin{center}
  {\Large\bfseries Wykład 2 --- 2026-10-09}
\end{center}

\section*{Nieprzywiedlność}

Odtąd, do odwołania, $(X_n)_{n \in \N_0}$ będzie jednorodnym łańcuchem Markowa
na co najwyżej przeliczalnej przestrzeni stanów $S$, o (jednokrokowych)
prawdopodobieństwach przejścia $p(i,j)$ (macierzowo:
$\mathbf{P} = \big(p(i,j)\big)_{i,j \in S}$) i rozkładzie początkowym $\alpha$.

\begin{definicja}[Osiągalność i komunikacja]
Niech $i,j \in S$. Powiemy, że
\begin{itemize}[topsep=0pt,itemsep=0pt]
  \item $j$ jest \emph{osiągalny} z $i$, jeśli
  $\exists n \in \N_0\colon\ p_n(i,j) > 0$ \ (oznaczenie: $i \to j$),
  \item $i,j$ \emph{komunikują się}, jeśli $i \to j$ i $j \to i$
  \ (oznaczenie: $i \leftrightarrow j$).
\end{itemize}
\end{definicja}

\begin{stwierdzenie}
Komunikacja jest relacją równoważności na zbiorze stanów $S$.
\end{stwierdzenie}

\begin{definicja}[Klasy komunikacji, nieprzywiedlność]
\begin{itemize}[topsep=0pt,itemsep=0pt]
  \item Klasy abstrakcji względem relacji komunikacji nazywają się
  \emph{klasami komunikacji} (ang.\ \emph{communication classes}).
  \item Łańcuch $(X_n)_n$ nazywamy \emph{nieprzywiedlnym}, ewentualnie
  \emph{nieredukowalnym} (ang.\ \emph{irreducible}), jeśli $i \leftrightarrow j$
  dla wszystkich $i,j \in S$; w przeciwnym razie mówimy, że $(X_n)_n$ jest
  \emph{przywiedlny} (ang.\ \emph{reducible}).
\end{itemize}
\end{definicja}

\section*{Okres stanu. Nieokresowość}

\begin{definicja}[Okres stanu]
\emph{Okresem} (ang.\ \emph{period}) stanu $j \in S$ nazywamy liczbę
\[
  o(j) := \nwd\{n \geq 1 : p_n(j,j) > 0\}.
\]
Jeśli stan $j$ nie jest osiągalny z samego siebie, wówczas $o(j)$ jest
największym wspólnym dzielnikiem zbioru pustego --- kładziemy wtedy
$o(j) = +\infty$. Jeżeli $o(i) = 1$, to stan $i$ nazywa się
\emph{nieokresowym} (ang.\ \emph{aperiodic}), a w przeciwnym razie
\emph{okresowym} (ang.\ \emph{periodic}).
\end{definicja}

\begin{zadanie}{(Minizadanie 1)}
\begin{enumerate}[label=\arabic*.]
  \item Zbadać okres każdego ze stanów łańcucha o macierzy przejścia
  \[
    \mathbf{P} =
    \begin{bmatrix}
      0 & \frac12 & \frac12 \\
      1 & 0 & 0 \\
      1 & 0 & 0
    \end{bmatrix}.
  \]
  Co się stanie z okresem każdego ze stanów, jeśli graf opisujący jednokrokową
  dynamikę łańcucha zmodyfikujemy poprzez dodanie pętli w stanie $1$
  (oczywiście taka zmiana musi pociągać zmianę w macierzy przejścia)?
  \item Podać przykład łańcucha Markowa (równoważnie, macierzy przejścia)
  o okresie $d \in \N$.
  \item Podać przykład możliwie małego (w sensie liczebności przestrzeni stanów)
  łańcucha, w którym istnieją zarówno stany o okresie $2$, jak i stany
  o okresie $3$.
  \item Zauważyć, że jeśli $p(j,j) = 1$, to $o(j) = 1$.
  \item Ile wynosi $o(j)$, jeśli istnieje $n$ taki, że zarówno $p_n(j,j) > 0$,
  jak i $p_{n+1}(j,j) > 0$?
\end{enumerate}
\end{zadanie}

Wartość okresu jest własnością klas komunikujących, tzn.\ zachodzi następujące

\begin{stwierdzenie}
Jeśli $i \leftrightarrow j$, to $o(i) = o(j)$.\footnote{Oczywiście implikacja
odwrotna nie jest prawdziwa --- dlaczego?}
\end{stwierdzenie}

\begin{wniosek}
W nieprzywiedlnym łańcuchu Markowa wszystkie stany mają ten sam okres.
\end{wniosek}

Jeśli więc wszystkie stany łańcucha nieprzywiedlnego są nieokresowe, łańcuch
nazywamy \textbf{nieokresowym}. Jeśli okresem każdego stanu w łańcuchu
nieprzywiedlnym jest liczba $d \neq 1$, łańcuch nazywamy \textbf{okresowym,
o okresie $d$}.\footnote{Podobnie można mówić o klasach okresowych o okresie
$d$.}

Można wykazać, że prawdziwy jest następujący opis struktury przestrzeni stanów
skończonych (łańcuch \textbf{skończony} to taki, którego przestrzeń stanów jest
zbiorem skończonym) okresowych łańcuchów Markowa:

\begin{stwierdzenie}
Niech $(X_n)_n$ --- nieprzywiedlny łańcuch Markowa na skończonej przestrzeni
stanów $S$. Jeśli $(X_n)_n$ jest okresowy, o okresie $d$, to istnieją parami
rozłączne $S_0, \ldots, S_{d-1} \subset S$ takie, że
$S = S_0 \cup \ldots \cup S_{d-1}$ i dla dowolnych $i,j \in S$ zachodzi
implikacja
\[
  p(i,j) > 0 \ \ \&\ \ i \in S_k \text{ dla pewnego } k \in \{0,\ldots,d-1\}
  \implies j \in S_l,
\]
gdzie $l = k+1 \mod d$.
\end{stwierdzenie}

Innymi słowy, przestrzeń stanów dzieli się na $d$ rozłącznych podzbiorów,
a przejścia w jednym kroku możliwe są tylko pomiędzy stanami należącymi do
dwóch kolejnych (modulo $d$) podzbiorów.

\begin{zadanie}{(Minizadanie 2)}
Opisać strukturę przestrzeni stanów nieprzywiedlnych i nieokresowych łańcuchów
Markowa z Minizadania 1.
\end{zadanie}

\section*{Chwilowość i powracalność}

Poniżej $(X_n)_n$ oznaczać będzie jednorodny łańcuch Markowa na co najwyżej
przeliczalnej przestrzeni stanów $S$. Jednokrokowe prawdopodobieństwa przejścia
oznaczamy dla $i,j \in S$ jako $p_{i,j}$ --- macierz przejścia to
$\mathbf{P} = [p_{i,j}]_{i,j \in S}$. Prawdopodobieństwa przejścia w $n$
krokach to $p_n(i,j)$. Rozkład początkowy reprezentujemy za pomocą
(wierszowego) wektora prawdopodobieństwa $\alpha = (\alpha_j)_{j \in S}$.

Dla $j \in S$ rozważamy moment $T_j$ pierwszej wizyty (po chwili $0$)
w stanie $j$
\[
  T_j = \min\{n \geq 1 : X_n = j\}
\]
(jeśli $X_n(\omega) \neq j$ dla wszystkich $n \in \N$, to kładziemy
$T_j(\omega) = \infty$). Zatem $T_j$ jest zmienną losową o wartościach
w $\N \cup \{\infty\}$.

\begin{uwaga}[Moment zatrzymania]
Zmienna $T_j$ jest przykładem \emph{momentu zatrzymania} lub \emph{momentu
stopu} (ang.\ \emph{stopping time}) względem tzw.\ \emph{filtracji naturalnej}
(ang.\ \emph{natural filtration}) łańcucha $(X_n)_n$, a więc niemalejącego
ciągu $\sigma$-ciał $\big(\sigma(X_0, \ldots, X_n)\big)_n$.

Ogólnie, jeśli $(\Omega, \mathcal{F}, \P)$ jest przestrzenią probabilistyczną,
\emph{filtracją} nazywamy dowolny niemalejący (w sensie inkluzji) ciąg
$(\mathcal{F}_n)_n$ pod-$\sigma$-ciał $\mathcal{F}$, czyli
$\mathcal{F}_n \subset \mathcal{F}$ dla dowolnego $n$, zaś
$\mathcal{F}_n \subset \mathcal{F}_m$ dla $n < m$.

Kiedy rozważamy proces stochastyczny (niekoniecznie łańcuch Markowa!)
określony na $(\Omega, \mathcal{F}, \P)$, to ma sens rozważać \emph{filtrację
naturalną procesu} $(X_n)_n$, a więc ciąg $(\mathcal{F}_n)_n$, zdefiniowany
\[
  \mathcal{F}_n := \sigma(X_0, \ldots, X_n).
\]
Takie pod-$\sigma$-ciało rozumiemy jako zawierające zdarzenia dotyczące procesu
$(X_n)_n$ zachodzące w chwili $n$, zaś filtracja $(\mathcal{F}_n)_n$ opisuje,
jak zmieniają się takie zdarzenia wraz z upływem czasu $n$.

Ogólnie, dla danej przestrzeni probabilistycznej $(\Omega, \mathcal{F}, \P)$
i filtracji $(\mathcal{F}_n)_n$, funkcja
$T \colon \Omega \to \N \cup \{+\infty\}$ jest \emph{momentem zatrzymania
względem} $(\mathcal{F}_n)_n$ (ewentualnie \emph{momentem stopu} lub
\emph{momentem Markowa}), jeśli
\[
  \{T = n\} \in \mathcal{F}_n
\]
dla każdego $n$. Chodzi więc o to, by wiedzieć w chwili $n$, czy
$T = n$.\footnote{Naszą wiedzę na temat procesu w chwili $n$ reprezentuje
$\mathcal{F}_n$.}

Jak wspomniano wyżej --- $T_j$ jest momentem stopu względem filtracji naturalnej
łańcucha. Oznacza to, że dla dowolnego $n \in \N$
\[
  \{T_j = n\} \in \sigma(X_0, \ldots, X_n).
\]
Tak z kolei jest, ponieważ
\[
  \{T_j = n\} = \{X_1 \neq j, \ldots, X_{n-1} \neq j, X_n = j\}.
\]
\end{uwaga}

\begin{uwaga}[Mocna własność Markowa]
Pojęcie momentu stopu związane jest ściśle z tzw.\ \emph{mocną własnością
Markowa} (ang.\ \emph{strong Markov property}). Jedno ze sformułowań tej
własności brzmi następująco: jeśli $(X_n)_n$ jest łańcuchem Markowa o macierzy
przejścia $\mathbf{P}$, a $T$ --- momentem stopu względem filtracji naturalnej,
to ciąg zmiennych losowych $(Y_n)_{n \in \N_0}$, dany przez
\[
  Y_n := X_{T+n}
\]
także jest łańcuchem Markowa, o takiej samej macierzy przejścia $\mathbf{P}$.
Intuicyjnie oznacza to tyle, że łańcuch w momencie stopu zapomina o przeszłości
i startuje od nowa od stanu $X_T$.
\end{uwaga}

Dla $i,j \in S$ oznaczmy przez $f_{ij}$ prawdopodobieństwo tego, że startując
z $i$ kiedykolwiek osiągniemy $j$, czyli
\[
  f_{ij}
  = \P\Big(\bigcup_{n=1}^{\infty} \{X_n = j\} \,\Big|\, X_0 = i\Big)
  = \P\Big(\bigcup_{n=1}^{\infty} \{T_j = n\} \,\Big|\, X_0 = i\Big)
  = \P\big(T_j < \infty \mid X_0 = i\big).
\]

Zdefiniujmy także zmienną losową zliczającą liczbę wizyt w stanie (poza chwilą
$n = 0$):\footnote{$N_j$ także jest zmienną losową o wartościach
w $\N \cup \{\infty\}$.}
\[
  N_j := \sum_{n=1}^{\infty} \ind_{\{X_n = j\}}.
\]

Nietrudno udowodnić, że dla dowolnych $i,j \in S$ i $k \in \N$
\begin{align*}
  \P\big(N_j \geq 1 \mid X_0 = i\big) &= f_{ij}, \\
  \P\big(N_j \geq k \mid X_0 = i\big) &= f_{ij} f_{jj}^{k-1}
\end{align*}
(w szczególności $\P\big(N_j \geq k \mid X_0 = j\big) = f_{jj}^k$).%
\footnote{Jak należy zinterpretować te równości?}
Oczywiście, powyższe równości wynikają z własności Markowa,\footnote{W jaki
sposób wynikają?} w szczególności do ich wykazania można posłużyć się
obserwacją (wynikającą z markowskości): jeśli $k < n$, a
$i_0, i_1, \ldots, i_n \in S$, to
\begin{align*}
  &\P\big(X_1 = i_1, \ldots, X_n = i_n \mid X_0 = i_0\big) \\
  &\quad= \P\big(X_1 = i_1, \ldots, X_k = i_k \mid X_0 = i_0\big)\,
    \P\big(X_{k+1} = i_{k+1}, \ldots, X_n = i_n \mid X_k = i_k\big) \\
  &\quad= \P\big(X_1 = i_1, \ldots, X_k = i_k \mid X_0 = i_0\big)\,
    \P\big(X_1 = i_{k+1}, \ldots, X_{n-k} = i_n \mid X_0 = i_k\big),
\end{align*}
w której ostatnia równość wynika ze stacjonarności (czyli niezmienniczości na
przesunięcia czasowe rozkładów warunkowych łańcucha Markowa) ---
jak?\footnote{Jak można interpretować tę równość?}

\begin{definicja}[Stan powracający i chwilowy]
Stan $i \in S$ nazywa się
\begin{itemize}[topsep=0pt,itemsep=0pt]
  \item \emph{powracającym}, ewentualnie \emph{powracalnym}
  (ang.\ \emph{recurrent}), jeśli $f_{ii} = 1$,
  \item \emph{chwilowym} (ang.\ \emph{transient}), jeśli $f_{ii} < 1$.
\end{itemize}
\end{definicja}

Zatem
\begin{itemize}[topsep=0pt]
  \item jeśli stan $j$ jest powracający, to
  $\P\big(N_j \geq k \mid X_0 = j\big) = 1$ dla każdego $k$
  (z prawdopodobieństwem $1$ łańcuch wystartowany z $j$ powraca do $j$
  nieskończenie wiele razy),
  \item jeśli stan $j$ jest chwilowy, to
  $\lim_{k \to \infty} \P\big(N_j \geq k \mid X_0 = i\big) = 0$ dla każdego
  $i \in S$, bez względu na jego chwilowość czy powracalność (na prawie każdej
  realizacji łańcucha wystartowanego z dowolnego $i$ stan $j$ jest odwiedzany
  skończenie wiele razy).
\end{itemize}

\begin{przyklad}
Jeśli macierzą przejścia łańcucha $(X_n)_n$ na $S = [3]$ jest
\[
  \mathbf{P} =
  \begin{bmatrix}
    \frac13 & \frac23 & 0 \\
    1 & 0 & 0 \\
    \frac14 & \frac12 & \frac14
  \end{bmatrix},
\]
to $f_{11} = 1$, $f_{22} = 1$ (bo szansa, że łańcuch na zawsze pozostanie w $1$
wynosi $\lim_n (1/3)^n = 0$), a $f_{33} < 1$, czyli $1$ i $2$ są powracające,
zaś $3$ jest chwilowy.
\end{przyklad}

Prawdziwe jest następujące twierdzenie:

\begin{twierdzenie}
Niech $j \in S$.
\begin{itemize}[topsep=0pt,itemsep=0pt]
  \item $j$ jest powracający
  \begin{align*}
    &\iff \P\big(N_j = \infty \mid X_0 = j\big) = 1 \\
    &\iff \sum_{n=1}^{\infty} p_n(j,j) = \infty,
  \end{align*}
  \item $j$ jest chwilowy
  \begin{align*}
    &\iff \P\big(N_j = \infty \mid X_0 = i\big) = 0 \quad \forall i \in S \\
    &\iff \sum_{n=1}^{\infty} p_n(i,j) = \frac{f_{ij}}{1 - f_{jj}} < \infty
      \quad \forall i \in S.
  \end{align*}
\end{itemize}
\end{twierdzenie}

Przy tym $\sum_{n=1}^{\infty} p_n(j,j) = \E\big[N_j \mid X_0 = j\big]$ ---
dlaczego?

Zauważmy, że z powyższego wynika, że jeśli $j$ jest chwilowy, to dla dowolnego
$i \in S$ zachodzi $\lim_{n \to \infty} p_n(i,j) = 0$.

\end{document}
